\question[4]
In Java wird bei einem Type Cast in einen kleineren ganzzahligen Typ einfach nur der Teil der Bits übernommen,
der in den Zieltyp passt. Ein \texttt{byte} hat 8 Bit und wird im Zweierkomplement gespeichert.
\begin{parts}
\part Bestimme den Wert von \texttt{(byte) 200}. Gib die Rechnung im Binärsystem an.
\part Bestimme den Wert von \texttt{(byte) 300}.
\part Beschreibe das Problem, das bei solchen Typumwandlungen auftritt.
\end{parts}

\begin{solutionbox}{6cm}
\begin{parts}
\part $200 = 11001000_2$. Das höchste Bit ist 1, also ist die Zahl im Zweierkomplement negativ:
$-128 + 64 + 8 = -56$.
\part $300 = 1\,00101100_2$ hat 9 Bit. Die unteren 8 Bit \texttt{00101100} werden übernommen: $32 + 8 + 4 = 44$.
\part Liegt der Wert außerhalb des Wertebereichs des Zieltyps, entsteht ohne Fehlermeldung ein ganz anderer Wert,
sogar mit anderem Vorzeichen.
\end{parts}
\end{solutionbox}
